% Master Diagram 3 -- Parity 1-bit nMOS cell
% Base for Q17 (CMOS version adds a p-block), Q18, Q20
% Depletion loads are shown as NMOS with gate tied to source (standard convention).
\documentclass[border=10pt]{standalone}
\usepackage{circuitikz}
\begin{document}
\begin{circuitikz}
% rails
\draw (-2,8) -- (5.5,8) node[right]{$V_{DD}$};
\draw (-2,0) -- (5.5,0);
\path (2,0) node[ground]{};
% depletion pull-up load (gate tied to source)
\path (0,6.4) node[nmos] (D1){};
\draw (D1.D) -- (0,8);
\draw (D1.S) -- (0,5.2) coordinate (PB);
\path (PB) node[circ]{};
\draw (D1.G) -| (-1.1,5.6) -- (0,5.6);
% Pbar_i node + bus
\draw (PB) -- (0,4.6);
\draw (-0.8,4.6) -- (1.0,4.6);
\path (1.0,4.7) node[above]{$\overline{P}_i$};
% branch 1
\path (-0.8,3.4) node[nmos] (T1a){}; \draw (T1a.D) -- (-0.8,4.6);
\draw (T1a.G) -- (-1.9,3.4) node[left]{$A_i$};
\path (-0.8,1.6) node[nmos] (T1b){}; \draw (T1b.D) -- (T1a.S);
\draw (T1b.G) -- (-1.9,1.6) node[left]{$\overline{P}_{i-1}$};
\draw (T1b.S) -- (-0.8,0);
% branch 2 (mirrored so gates face right)
\path (0.8,3.4) node[nmos, xscale=-1] (T2a){}; \draw (T2a.D) -- (0.8,4.6);
\draw (T2a.G) -- (1.9,3.4) node[right]{$\overline{A}_i$};
\path (0.8,1.6) node[nmos, xscale=-1] (T2b){}; \draw (T2b.D) -- (T2a.S);
\draw (T2b.G) -- (1.9,1.6) node[right]{$P_{i-1}$};
\draw (T2b.S) -- (0.8,0);
% static inverter
\path (4,6.4) node[nmos] (D2){};
\draw (D2.D) -- (4,8);
\draw (D2.S) -- (4,5.2) coordinate (PO);
\path (PO) node[circ]{};
\draw (D2.G) -| (3.0,5.6) -- (4,5.6);
\path (4,2.8) node[nmos] (T3){};
\draw (T3.D) -- (PO);
\draw (T3.S) -- (4,0);
% Pbar_i drives inverter input
\draw (0,4.6) -- (2.3,4.6) -- (2.3,2.8) -- (T3.G);
% output
\draw (PO) -- (5.3,5.2) node[right]{$P_i$};
\end{circuitikz}
\end{document}
